\documentclass[11pt]{article}
\usepackage[T1]{fontenc}
\usepackage{lmodern,amsmath,amssymb,amsthm,mathtools,microtype}
\usepackage[margin=27mm]{geometry}
\usepackage[colorlinks=true,linkcolor=blue,citecolor=blue,urlcolor=blue]{hyperref}
\setlength{\emergencystretch}{2em}
\newcommand{\Ric}{\operatorname{Ric}}
\newcommand{\grad}{\operatorname{grad}}
\newcommand{\tr}{\operatorname{tr}}
\newcommand{\diver}{\operatorname{div}}
\newcommand{\dd}{\,dV_g}
\newtheorem{theorem}{Theorem}
\newtheorem{corollary}[theorem]{Corollary}
\newtheorem{proposition}[theorem]{Proposition}
\theoremstyle{remark}
\newtheorem*{remark}{Remark}
\title{Traceless-Ricci Descent and Harmonic-Flux Equality\\
\large A continuation of the smooth global minimizer problem}
\author{}
\date{8 October 2026}
\begin{document}
\maketitle
\begin{abstract}
For a closed almost-K\"ahler four-manifold with $J$-invariant Ricci tensor,
we derive an exact first variation of the traceless-Ricci energy from the
classical Bach identity of Apostolov--Dr\u aghici. It gives an explicit
volume-preserving descending deformation in the nonpositive-scalar range.
Combined with the supplied manuscript's harmonic-flux equality theorem,
it shows that no non-Bach-flat critical metric can lie in that equality
class. This eliminates a candidate construction route for the smooth
global minimizer problem; it does not establish smooth attainment.
\end{abstract}

\section{Convention and the classical input}

All manifolds below are smooth, closed, connected and four-dimensional.
For a Riemannian metric set
\[
 E=\Ric-\frac{s}{4}g,\quad
 P(g)=\int_M|E|^2\dd,\quad S(g)=\int_Ms^2\dd,\quad
 \mathcal W(g)=\int_M|W_{\mathrm{op}}|^2\dd.
\]
The Weyl norm is the curvature-operator norm. Chern--Gauss--Bonnet reads
\begin{equation}
 P=2\mathcal W+\frac{S}{12}-16\pi^2\chi(M).
 \label{eq:GB}
\end{equation}
Use the Bach normalization $B=\grad\mathcal W$ of
\cite[\S2.3, equation (15)]{AD}. This specifies its sign and coefficient;
vanishing is independent of that choice. For an almost-K\"ahler structure
$(g,J,\Omega)$ with $J$-invariant Ricci, their Lemma~3 gives
\begin{equation}
 B=-\frac13(\nabla ds)^{\mathrm{inv}}
   +\frac16(\nabla ds)^{\mathrm{anti}}
   -\frac{\Delta s}{12}g-\frac{s}{6}E+\mathcal T,
 \label{eq:AD}
\end{equation}
where $\mathcal T$ is symmetric and $J$-anti-invariant, and
$\Delta=d^*d$ is nonnegative. We use this established identity, rather
than claim a new formula for the Bach tensor.

The contracted Bianchi identity gives $\diver E=ds/4$. Since $E$ is
trace-free and $J$-invariant, integrating \eqref{eq:AD} against $E$ yields
\begin{equation}
 D\mathcal W_g[E]
 =\int_M\langle B,E\rangle\dd
 =\frac1{12}\int_M|ds|^2\dd-\frac16\int_Ms|E|^2\dd.
 \label{eq:B-pair}
\end{equation}
Indeed, the anti-invariant terms are orthogonal to $E$, and
$\int\langle\nabla ds,E\rangle=-\int|ds|^2/4$.

\section{A genuine global metric deformation}

\begin{theorem}\label{thm:descent}
Let $(M,g,J,\Omega)$ be closed almost-K\"ahler, and suppose
$\Ric(JX,JY)=\Ric(X,Y)$. Then
\begin{equation}
 \boxed{DP_g[E]=\frac18\int_M|ds|^2\dd
                    -\frac12\int_Ms|E|^2\dd.}
 \label{eq:main}
\end{equation}
The smooth global metric path
\begin{equation}
 g_t(X,Y)=g\bigl(\exp(-tE^\sharp)X,Y\bigr)
 \label{eq:path}
\end{equation}
is positive definite and preserves the volume form for every real $t$.
If $s\le0$, it strictly decreases $P$ for all sufficiently small $t>0$,
unless $ds\equiv0$ and $sE\equiv0$.
\end{theorem}

\begin{proof}
For a covariant metric variation $h$,
\[
 s'[h]=-\langle\Ric,h\rangle
       +\diver\diver h-\Delta_{\mathrm{div}}\tr h,
 \qquad \Delta_{\mathrm{div}}=\diver\nabla=-\Delta.
\]
For $h=E$, this becomes
$s'[E]=-|E|^2+\Delta_{\mathrm{div}}s/4$, while
$dV'[E]=0$. Thus
\begin{equation}
 DS_g[E]=-2\int_Ms|E|^2\dd-\frac12\int_M|ds|^2\dd.
 \label{eq:S-pair}
\end{equation}
Combine \eqref{eq:GB}, \eqref{eq:B-pair} and \eqref{eq:S-pair} to obtain
\eqref{eq:main}. The endomorphism $E^\sharp$ is self-adjoint and
trace-free. Consequently $\exp(-tE^\sharp)$ is positive definite with
determinant one. Hence \eqref{eq:path} is an actual metric path, not
a formal curvature variation, and has fixed volume form. Its derivative
at zero is $-E$, giving
\begin{equation}
 \left.\frac{d}{dt}P(g_t)\right|_{t=0}
 =-\frac18\int_M|ds|^2\dd+\frac12\int_Ms|E|^2\dd.
 \label{eq:negative}
\end{equation}
Both terms are nonpositive when $s\le0$. Equality requires both
$ds=0$ and $s|E|^2=0$; connectedness gives the stated exceptions.
\end{proof}

The varied metrics need not remain almost-K\"ahler. This is appropriate
for an infimum over \emph{all} smooth metrics, including a unit-volume
normalization. In particular the deformation is not a claim about a
minimum restricted to a fixed symplectic form.

\begin{corollary}\label{cor:csc}
For a constant-scalar metric in Theorem~\ref{thm:descent},
\[
 DP_g[E]=-\frac{s}{2}P(g).
\]
If $s\ne0$ and $P(g)>0$, the path
$g(\exp(t\operatorname{sgn}(s)E^\sharp)\cdot,\cdot)$ has derivative
$-|s|P(g)/2<0$. Thus a constant-nonzero-scalar $P$-critical metric in
this entire almost-K\"ahler, $J$-invariant-Ricci class is Einstein.
\end{corollary}

If $s<0$, one may also conformally correct this path to retain the initial
constant scalar curvature. The linearization of the equation
$6\Delta_g u+s_gu=s_0u^3$ at $u=1$, $s_g=s_0<0$, is
$6\Delta_g-2s_0$, which is invertible. The implicit function theorem
gives the correction. Since the conformal first variation of $P$ vanishes
at a constant-scalar metric, the strict derivative is unchanged.

\begin{corollary}
If $s\le-\sigma<0$ and $P>0$ in Theorem~\ref{thm:descent}, then
\[
 \|\grad P\|_{L^2}
 \ge \frac{\|ds\|_{L^2}^2}{8\sqrt P}
       +\frac{\sigma}{2}\sqrt P.
\]
\end{corollary}
\begin{proof}
Apply Cauchy--Schwarz to $DP[E]=\int\langle\grad P,E\rangle$ and
use \eqref{eq:main} and $\|E\|_2=\sqrt P$.
\end{proof}
This estimate concerns the exact geometric class in the theorem.
It does not supply the derivative control needed to extend the identity
to arbitrary nearly aligned, degenerating metrics.

\section{Consequence for the supplied harmonic-flux decomposition}

Write $\mathcal B:\Lambda^+\to\Lambda^-$ for the off-diagonal curvature
block, so $P=4\|\mathcal B\|_2^2$. For a nonzero harmonic anti-self-dual
two-form $h$, put
\[
 H=\|h\|_2^2,\qquad F_h=\sqrt2\int_M|\mathcal B^*h|\dd,
 \qquad L_h=\frac{P/2-F_h^2/H}{4\pi^2}\ge0.
\]
The supplied continuation \cite[Sections 4--5]{Input} proves the following
equality statement, including removal across the harmonic zero set:
\begin{quote}
At constant scalar curvature, $L_h=0$ implies that $g$ is Einstein or
admits a smooth global almost-K\"ahler structure with $J$-invariant Ricci.
\end{quote}
This is the precise additional input used in the next proposition.

\begin{proposition}\label{prop:flux}
Assuming the equality statement above, a smooth closed $P$-critical metric
with $L_h=0$ for some nonzero harmonic anti-self-dual $h$ is Bach-flat.
If its scalar curvature is nonzero, it is Einstein.
\end{proposition}
\begin{proof}
For $g_t=e^{2tf}g$, conformal invariance of $\mathcal W$ and
\eqref{eq:GB} give
\[
 DP_g[2fg]=\int_M f\Delta s\dd.
\]
Criticality implies that $s$ is constant. If $s\ne0$, the equality
statement and Corollary~\ref{cor:csc} force $E=0$. If $s=0$, the
first variation of $S$ vanishes, and \eqref{eq:GB} gives
$\grad P=2B$. Criticality then implies $B=0$.
\end{proof}

Therefore, on the negative-Yamabe candidates with $P_*>0$ already
identified in the project, a hypothetical smooth global minimizer must have
\begin{equation}
 L_h(g)>0\qquad\text{for every nonzero }h\in\mathcal H_g^-.
 \label{eq:positive-remainder}
\end{equation}
At a fixed such metric with $b_->0$, compactness of the harmonic unit
sphere and continuity give $\min_{\|h\|_2=1}L_h(g)>0$.
This last positive number depends on the metric. It is not a uniform
bound along a minimizing sequence.

The direction of this consequence matters: forcing $L_h=0$ would
\emph{exclude} the desired non-Bach-flat minimizer. It would not establish
its existence. A sharp inequality attained by such a minimizer must
accommodate a nonzero harmonic-alignment remainder or a different
geometric structure.

\section{Two normalization checks and the remaining problem}

For a product of constant-curvature surfaces with Gaussian curvatures
$a,b$ and product volume $V$, one has $s=2(a+b)$ and
$P=(a-b)^2V$. The eigenvalues of $E^\sharp$ on the two factors are
$(a-b)/2$ and $-(a-b)/2$. Under
$g_t=g(\exp(tE^\sharp)\cdot,\cdot)$,
\[
 P(g_t)=V\bigl(ae^{-t(a-b)/2}-be^{t(a-b)/2}\bigr)^2.
\]
Its derivative is $-(a+b)(a-b)^2V=-sP/2$, independently checking the
coefficient and sign of Corollary~\ref{cor:csc}.

There is also a direct distinction between the two global variational
problems in the supplied manuscript. Let
\[
 R=\int|\Ric|^2=P+S/4,\quad
 T_s=\int s,\quad V=\operatorname{Vol}(M),\quad
 \mathcal E=4-\frac{T_s^2}{VR}.
\]
At a smooth $P$-critical metric with nonzero constant scalar curvature,
the first variation of $S-T_s^2/V$ vanishes. Thus
\begin{equation}
 \boxed{D\mathcal E_g[E]=\frac{2sP^2}{R^2}.}
 \label{eq:distinct}
\end{equation}
Indeed, $\mathcal E=4P/R$ to first order there, $DP=0$, and
$DS[E]=-2sP$. If $P>0$, \eqref{eq:distinct} is nonzero. Hence smooth
attainment of the normalized defect does not produce the sought
non-Bach-flat $P$-minimizer.

\paragraph{What remains.}
The target is still a smooth global minimizer of $P$ on a closed manifold
which is not Bach-flat. No smooth-attainment theorem or such metric is
constructed here. The result proved here is the explicit descending
deformation and the consequent exclusion of the supplied flux-equality
class. The Bach identity is classical; the connection to the project's
flux equality is the continuation recorded in this note.

\begin{thebibliography}{9}
\bibitem{AD}
V. Apostolov and T. Dr\u aghici,
\emph{Almost K\"ahler 4-manifolds with $J$-invariant Ricci tensor and
special Weyl tensor}, Quart. J. Math. 51 (2000), 275--294.
\href{https://arxiv.org/abs/math/9911196}{arXiv:math/9911196}.
The input is \S2.3, equation (15), Lemma~3.
\bibitem{Input}
Supplied project manuscript, \emph{Einstein Gaps on Four-Manifolds:
Subsequent Mathematical Developments}, 8 October 2026,
\texttt{03\_Subsequent\_Mathematical\_Developments.tex},
Sections 4--5 (harmonic flux, equality and removability).
\end{thebibliography}
\end{document}
